\documentclass[11pt]{article}
\usepackage[margin=1in]{geometry}
\usepackage{enumitem}
\usepackage[T1]{fontenc}
\usepackage[utf8]{inputenc}

\begin{document}

\section*{PBI-001 Shariah Review Response Draft}

\textbf{Document status:} Draft working paper for Shariah review.\\
\textbf{Purpose:} Provide a standards-aligned recommendation for the MVP baseline before implementation PBIs are created.\\
\textbf{Important note:} This document is not a formal Shariah pronouncement. Final approval remains subject to the designated Shariah reviewer or Shariah Board.

\section*{Recommended Meeting Outcome}

\begin{itemize}[leftmargin=1.5em]
    \item \textbf{Decision / recommendation:} Approved with conditions.
    \item \textbf{Rationale:} The proposed single-venture restricted mudarabah is the strongest MVP candidate under the 12-week constraint because it is narrower, easier to govern, and easier to validate than a broader PLS design.
    \item \textbf{Conditions / caveats:}
    \begin{itemize}[leftmargin=1.5em]
        \item The contract must state a clear profit-sharing ratio and must not imply a guaranteed principal or guaranteed return.
        \item The restriction scope must be defined precisely enough to identify the permitted venture and use of funds.
        \item The deductible-cost taxonomy must be expressly approved before the first distribution event.
        \item Any negligence, misconduct, or breach determination must follow documented evidence review and must not be inferred automatically from a commercial loss alone.
        \item Any security, guarantee, or collateral structure must be reviewed separately to ensure it does not convert the arrangement into an impermissible principal or profit guarantee.
        \item Musharakah may be deferred for MVP, but the ADR and roadmap should state clearly that deferral is a sequencing choice, not a rejection of musharakah as a valid future model.
    \end{itemize}
    \item \textbf{Approval status:} Approved with conditions
    \item \textbf{Reviewer name and role:} \rule{2.5in}{0.4pt}
    \item \textbf{Date:} \rule{1.5in}{0.4pt}
\end{itemize}

\section*{Responses to Review Questions}

\begin{enumerate}[leftmargin=1.8em,label=\textbf{Q\arabic*.}]
    \item \textbf{Is the proposed ``single-venture restricted mudarabah'' acceptable as an MVP framing?}

    \textbf{Recommended answer:} Yes, subject to conditions. It is acceptable as an MVP framing because it preserves the core mudarabah structure while keeping scope narrow enough for governance, documentation, and implementation control. The contract should be framed as a restricted mudarabah tied to a specifically identified venture, with no guarantee of principal or profit.

    \item \textbf{Can the procurement reference (PO / awarded tender / job) be used as the restriction scope?}

    \textbf{Recommended answer:} Yes, in principle, provided the procurement reference is sufficiently specific to operate as the investment mandate. The reference should identify the venture, permitted use of funds, relevant parties, expected business activity, and evidence trail. A procurement reference may therefore be used as the operational restriction scope \emph{if} it is not merely a label, but a defined and auditable restriction boundary.

    \item \textbf{How should ``realized profit'' be defined in this context?}

    \textbf{Recommended answer:} For MVP purposes, ``realized profit'' should mean profit attributable to the restricted venture for the relevant distribution period, calculated using the approved profit methodology, supported by documentary evidence, and not based on a target rate, estimate presented as certainty, or any guaranteed minimum return. A practical formulation is:

    \begin{quote}
    Realized profit = recognized venture revenue or value realized for the approved period,
    minus approved direct venture costs, minus recognized provisions or losses attributable
    to that venture, as determined under the approved calculation basis.
    \end{quote}

    Realized profit should be distributed only after the agreed calculation basis and deductible-cost policy have been approved.

    \item \textbf{Which cost categories may be deducted before profit distribution?}

    \textbf{Recommended answer:} Use a conservative two-tier policy.

    \textbf{Deductible before distributable profit:}
    \begin{itemize}[leftmargin=1.5em]
        \item direct procurement or venture acquisition cost;
        \item direct delivery, fulfilment, or logistics cost attributable to the restricted venture;
        \item direct transaction fees, taxes, or charges tied to that venture;
        \item directly attributable impairment, write-down, or provision related to that venture;
        \item direct recovery, liquidation, or enforcement cost where contractually relevant.
    \end{itemize}

    \textbf{Not deductible before distributable profit unless separately disclosed and expressly approved:}
    \begin{itemize}[leftmargin=1.5em]
        \item general platform overhead;
        \item unrelated salaries and corporate administrative expenses;
        \item broad technology or hosting cost;
        \item business-development cost;
        \item discretionary buffers or smoothing reserves;
        \item any amount that functions in substance like a guaranteed return or hidden management charge.
    \end{itemize}

    \textbf{Implementation note:} The deductible-cost policy should be versioned and approved as a named rule set before any distribution.

    \item \textbf{Is manual approval required before every distribution event?}

    \textbf{Recommended answer:} Not necessarily before \emph{every} event, provided the product structure, profit formula, deductible-cost taxonomy, and evidence requirements have already been approved and the system performs compliance checks and records review evidence. However, manual review should be mandatory for:
    \begin{itemize}[leftmargin=1.5em]
        \item the first live distribution under a new template or rule version;
        \item any exception or override;
        \item any change to the profit formula or deductible-cost mapping;
        \item any event involving loss reallocation to the mudharib;
        \item any suspected Shariah non-compliance or dispute.
    \end{itemize}

    \item \textbf{What evidence is sufficient to determine negligence / misconduct / breach?}

    \textbf{Recommended answer:} Commercial loss alone is not sufficient. For MVP purposes, a negligence, misconduct, or breach finding should require a documented evidence pack showing both:
    \begin{itemize}[leftmargin=1.5em]
        \item the applicable contractual restriction or duty; and
        \item objective evidence that the mudharib acted outside that restriction or failed that duty.
    \end{itemize}

    \textbf{Minimum evidence pack:}
    \begin{itemize}[leftmargin=1.5em]
        \item executed contract and approved restriction scope;
        \item procurement reference and supporting commercial documents;
        \item disbursement records and designated-account trail;
        \item invoices, purchase orders, goods receipt, delivery evidence, or service completion evidence;
        \item exception logs, approvals, and amendment history;
        \item internal review memo describing the alleged negligence, misconduct, or breach;
        \item supporting correspondence or incident evidence;
        \item reviewer conclusion and approval record.
    \end{itemize}

    \textbf{Implementation note:} No automated rule should reallocate loss to the mudharib without documented human review.

    \item \textbf{Are any security or collateral arrangements permissible without becoming an impermissible guarantee?}

    \textbf{Recommended answer:} Yes, but only with careful structuring. The safest MVP position is:
    \begin{itemize}[leftmargin=1.5em]
        \item do \emph{not} allow the mudharib to guarantee principal or profit as part of the mudarabah bargain;
        \item allow review of third-party guarantee structures only if they are separately justified and approved;
        \item allow collateral or security only where it secures a permissible claim, such as misconduct, negligence, breach, or an independently valid guarantee structure;
        \item do not let collateral wording imply that ordinary commercial loss will automatically be repaid.
    \end{itemize}

    \textbf{MVP policy recommendation:} Permit only pre-approved security patterns, and require explicit Shariah review of any guarantee or collateral language before use.

    \item \textbf{Are there any mandatory disclosure or documentation fields missing from the draft schema?}

    \textbf{Recommended answer:} Yes. In addition to the draft fields, add the following:
    \begin{itemize}[leftmargin=1.5em]
        \item \texttt{venture\_description};
        \item \texttt{permitted\_use\_of\_funds};
        \item \texttt{shariah\_pronouncement\_reference};
        \item \texttt{distribution\_period\_start} and \texttt{distribution\_period\_end};
        \item \texttt{approved\_profit\_definition};
        \item \texttt{deductible\_cost\_policy\_version};
        \item \texttt{designated\_collection\_account} and \texttt{designated\_disbursement\_account};
        \item \texttt{evidence\_bundle\_id};
        \item \texttt{security\_or\_collateral\_details}, if any;
        \item \texttt{recourse\_on\_breach};
        \item \texttt{distribution\_approval\_status};
        \item \texttt{exception\_flag};
        \item \texttt{amendment\_log};
        \item \texttt{shariah\_noncompliance\_flag} and \texttt{purification\_record}, if applicable.
    \end{itemize}

    \item \textbf{Is it acceptable to defer musharakah support to post-MVP?}

    \textbf{Recommended answer:} Yes. Deferring musharakah is acceptable as an MVP sequencing decision, provided the documentation states clearly that musharakah remains a valid future extension and is not being rejected on Shariah grounds. The MVP may adopt restricted mudarabah as the initial baseline while preserving a roadmap for later musharakah support.

    \item \textbf{What wording should appear in user-facing contract summaries to avoid implying guaranteed return?}

    \textbf{Recommended answer:} Use wording that is explicit, simple, and non-promissory. Suggested text:

    \begin{quote}
    This arrangement is a restricted mudarabah linked to Procurement Reference [X]. Profit, if any, will be shared between the capital provider and the venture operator according to the agreed profit-sharing ratio and the approved profit-calculation method. Returns are variable and depend on the actual performance of the venture. Principal and profit are not guaranteed. Financial loss is borne in accordance with mudarabah principles, except where negligence, misconduct, or breach by the mudharib is established.
    \end{quote}

    \textbf{Avoid the following terms in user-facing summaries unless separately qualified and approved:}
    \begin{itemize}[leftmargin=1.5em]
        \item guaranteed return;
        \item protected capital;
        \item fixed yield;
        \item assured payout;
        \item minimum profit;
        \item principal protection.
    \end{itemize}
\end{enumerate}

\section*{Primary Sources for Reviewer Pack}

\begin{enumerate}[leftmargin=1.8em]
    \item Securities Commission Malaysia, \emph{Guidelines on Islamic Capital Market Products and Services}, especially Appendix 7 definitions of \emph{mudharabah}, \emph{musharakah}, \emph{kafalah}, and \emph{rahn}.
    \item Securities Commission Malaysia, \emph{Resolutions of the Securities Commission Shariah Advisory Council}, especially:
    \begin{itemize}[leftmargin=1.5em]
        \item \emph{Kafalah on Mudharabah Capital};
        \item \emph{Ujrah for Guarantees}.
    \end{itemize}
    \item IFSB, \emph{Guiding Principles on Sharī`ah Governance Systems for Institutions Offering Islamic Financial Services} (IFSB-10).
    \item IFSB, \emph{FAQs on IFSB-10}, especially ex-ante approval, compliance checks, internal Shariah review, and Shariah governance reporting.
    \item IFSB, \emph{Revised Compilation Guide on Prudential and Structural Islamic Financial Indicators (PSIFIs)}, especially sections on:
    \begin{itemize}[leftmargin=1.5em]
        \item explicit contracted profit-sharing ratios;
        \item actual realized returns;
        \item distributions to investment account holders;
        \item distinction between direct financing/investment expenses and general operating expenses.
    \end{itemize}
\end{enumerate}

\section*{Supplementary Official Malaysian Sources for Direct Board Review}

\begin{enumerate}[leftmargin=1.8em]
    \item Bank Negara Malaysia, \emph{Shariah Standard on Mudarabah}.
    \item Bank Negara Malaysia, \emph{Shariah Governance Policy Document} (2019).
    \item Bank Negara Malaysia, \emph{Shariah Resolutions in Islamic Finance} (2nd edition).
\end{enumerate}

\section*{Suggested Meeting Record Template}

\begin{itemize}[leftmargin=1.5em]
    \item \textbf{Decision / recommendation:} Approved / Approved with conditions / Requires revision / Rejected
    \item \textbf{Conditions or caveats:} \rule{4.5in}{0.4pt}
    \item \textbf{Reviewer name and role:} \rule{3.5in}{0.4pt}
    \item \textbf{Date:} \rule{1.5in}{0.4pt}
    \item \textbf{Approval status:} \rule{2in}{0.4pt}
\end{itemize}

\end{document}